\chapter{\introductionname}
\label{chap:introduction}

\section[Context and Motivations]{Context and Motivations of the Thesis Project}

In modern enterprises, data is often viewed as a resource that can flow seamlessly from operational systems to analytical platforms. 
In practice, however, this perspective overlooks a fundamental issue: data is collected to support day-to-day business operations but is later expected to drive strategic decision-making. 
As a result, a structural misalignment emerges between \textbf{data generation} ("writing") and \textbf{data consumption for decision-making} ("reading").
    
This misalignment is also reflected at the organizational level. 
On the one hand, the Chief Information Officer (CIO) is responsible for systems designed to ensure operational efficiency, speed, and transactional integrity. 
On the other hand, the Chief Data Officer (CDO) is tasked with transforming the same data into meaningful business knowledge. 
As highlighted in the literature, data is not an objective representation of organizational reality; rather, it reflects the design choices, operational constraints, and organizational context in which it is generated \cite{boyd2012critical, kitchin2014data}.

When data is extracted from operational systems to support analytical tools or predictive models, it undergoes what Leonelli defines as a \textbf{data journey} \cite{Leonelli2016data}. 
Throughout this process, the different objectives of operational and analytical systems become evident. 
Operational systems, such as ERP, MES, and CRM, are designed to record individual transactions quickly and reliably. 
Analytical systems, instead, require semantically consistent data that can be compared over time and used to identify trends, patterns, and causal relationships. 
Consequently, operational data cannot be directly reused for analytical purposes without significant transformation \cite{borgman2015big}.

The literature also emphasizes that data classification structures and database schemas determine in advance which information can be captured and which remains invisible to the organization \cite{bowker1999sorting}. 
Therefore, data quality should not be treated solely as a downstream issue to be addressed through cleaning and integration processes. 
Instead, it should be considered from the moment data is generated. 
This requires a governance framework in which operational and analytical systems share a common semantic and structural foundation.

These concepts are explored through the case study of Amadori S.p.A., a leading Italian agri-food company specializing in a vertically integrated poultry supply chain. 
In an industry characterized by narrow profit margins, strict regulatory requirements, and highly perishable products, 
the ability to transform operational data into strategic information is essential for maintaining competitiveness.

The poultry supply chain makes this challenge even more complex because of the biological constraints that govern livestock production. 
Breeding, hatching, and growth cycles follow fixed timelines that cannot be adjusted in response to rapid changes in market demand. 
As a result, the production system is inherently rigid, while the downstream market is characterized by volatile demand, promotional campaigns, and strong seasonal fluctuations.

This challenge provides the motivation for the present thesis, which aims to develop an approach to reduce the gap between the transactional data generated by operational systems 
and the semantic information required by business decision-makers. 
The objective is to improve the integration between operational and analytical systems, 
making data more understandable and actionable in decision-making processes.


\section[Objective of the Work]{Objective of the Work: Innovating the Data Consumption Experience}

To address the misalignment between operational and analytical data in the agri-food supply chain, this thesis proposes the design and evaluation of an architectural framework centered on the integration of a centralized \textbf{Semantic Layer} and a \textbf{well-governed Data Catalog}.

The proposed approach differs from traditional Business Intelligence architectures, where new analytical requirements often require extensive data engineering efforts and modifications to database schemas. 
Instead, the proposed architecture introduces a \textbf{Semantic Layer} that abstracts heterogeneous operational systems and exposes a shared business representation of enterprise data.

While the Semantic Layer provides a unified business representation of enterprise data, the Data Catalog serves as the authoritative repository for metadata, 
business definitions, data lineage, ownership, and documentation. Together, these components establish a consistent semantic foundation, improve data discoverability, 
and facilitate the integration of analytical applications and AI-based services.

Building upon this foundation, the thesis explores the integration of \textbf{Large Language Model (LLM) agents} as a natural language interface to enterprise data. 
By leveraging the business semantics provided by the Semantic Layer together with the metadata managed by the Data Catalog, the agents can interpret users' requests and translate them into appropriate analytical queries.

This approach enables corporate planners to access operational insights without requiring knowledge of query languages or relying exclusively on predefined dashboards. 
Furthermore, because the AI Agents operate within the semantic constraints defined by the Semantic Layer and supported by the governed metadata of the Data Catalog, 
the risk of inconsistent interpretations and hallucinations is significantly reduced.

Ultimately, the proposed framework aims not only to improve data accessibility for business users but also to strengthen the role of the \textbf{Data Office}. 
Rather than focusing primarily on the production of reports and ad hoc data preparation, the Data Office evolves toward a governance function responsible for defining, 
maintaining, and certifying the semantic assets and metadata that enable trusted AI-driven analytics and decision support across the organization.

\section{Structure of the Thesis}

To systematically address the research objectives and the deployment of the proposed semantic framework within Amadori S.p.A., this thesis is organized into six interconnected chapters, 
transitioning from foundational organizational theory to technical implementation and practical validation:

\begin{itemize}
    \item \Cref{chap:context} focus on the poultry sector's unique dynamics and complexity, with a particular focus on the Demand Forecasting process.
    
    \item \Cref{chap:demandforecast} examines the quantitative foundation of the enterprise's predictive capabilities. It provides an overview of statistical forecasting models and diagnoses the limitations of the company's ``as-is'' \textit{ARIMAX} approach.
    
    \item \Cref{chap:architecture} defines the socio-technical framework implemented to overcome data silos. It details the paradigm shift toward \textit{Data Products}, the establishment of a \textit{Semantic Layer and Business Glossary}, and the strategic use of \textit{Ontologies and Knowledge Graphs}. Furthermore, it emphasizes the centrality of \textit{Metadata Management} and introduces the human factor via \textit{Data \& AI Champions}.
   
    \item \Cref{chap:techstack} outlines the concrete IT infrastructure enabling the transformation. It formalizes business data through Looker as the Semantic Layer and Interface Layer, introduces the OpenMetadata ecosystem for governance and ontology management, and details the component orchestration required to enable Generative AI.
    
    \item \Cref{chap:casestudy} focuses on translating demand forecasts into net raw material requirements within a simplified representation of the real business scenario. The chapter concludes by presenting the implementation of the use case, demonstrating how the Semantic Layer enables intuitive and semantically consistent analytical exploration.
    
    \item \Cref{chap:chatbot} represents the innovative climax of the thesis. It demonstrates how corporate ontologies mitigate LLM hallucinations by providing ``certified'' context, and provides a comprehensive simulation of the user experience from natural language queries to final operational decisions.
    
    \item \Cref{chap:conclusions} synthesizes the tangible business benefits achieved while outlining future vectors for scalability within the organization.
\end{itemize}
